% --- Archon physics formalization source begin ---
% archon:physics
% archon:covers PhyXMiniProblems/problem_phyx_mini_0775.lean
% archon:source-report reports/phyx_mini/problem_phyx_mini_0775.source.json

\chapter{Physics problem phyx\_mini\_0775}
\label{ch:PhyXMiniProblems_problem_phyx_mini_0775}

\paragraph{Problem source.}
\#\# Physical scenario

Two metal disks, one with radius \textbackslash{}(R\_1 = 2.50\textbackslash{}text\{ cm\}\textbackslash{}) and mass \textbackslash{}(M\_1 =0.80\textbackslash{}text\{ kg\}\textbackslash{}) and the other with radius \textbackslash{}(R\_2 = 5.00\textbackslash{}text\{ cm\}\textbackslash{}) and mass \textbackslash{}(M\_2 = 1.60\textbackslash{}text\{ kg\}\textbackslash{}), are welded together and mounted on a frictionless axis through their common center as shown in figure.

\#\# Question

A light string is wrapped around the edge of the smaller disk, and a \textbackslash{}(1.50\textbackslash{}text\{ kg\}\textbackslash{}) block is suspended from the free end of the string. If the block is released from rest at a distance of \textbackslash{}(2.00\textbackslash{}text\{ m\}\textbackslash{}) above the floor, what is its speed just before it strikes the floor?

\#\# Answer choices

- A: A: 3.12m/s

- B: B: 3.40m/s

- C: C: 3.86m.s

- D: D: 4.86m/s

\#\# Figure caption (auxiliary; use the image as primary evidence)

The image depicts a pulley system. The main components and details are:

1. **Pulley System:** 
   - There are two concentric pulleys with different radii indicated as \textbackslash{}( R\_1 \textbackslash{}) and \textbackslash{}( R\_2 \textbackslash{}).
   - The axis of the pulleys is horizontal.

2. **Radii:** 
   - \textbackslash{}( R\_1 \textbackslash{}) is the radius of the smaller pulley.
   - \textbackslash{}( R\_2 \textbackslash{}) is the radius of the larger pulley.

3. **Rope and Mass:**
   - A rope runs over one of the pulleys and is attached to a hanging mass.
   - The mass is labeled as 1.50 kg and is suspended directly below.

4. **Mass:** 
   - The mass is a solid block shape and is colored orange.

This setup likely represents part of a physics problem involving rotational dynamics or torque.

\#\# Dataset metadata

Rotational Motion; Physical Model Grounding Reasoning, Multi-Formula Reasoning

\paragraph{Recorded answer/context.}
B: B: 3.40m/s

\paragraph{Figure/image path.}
/root/proposal\_for\_physic/science-mango/phyx\_mini\_run/phyx\_data/test\_image/775.png

\paragraph{Formalization target.}
create a compiling Lean file with sorry bodies at `PhyXMiniProblems/problem\_phyx\_mini\_0775.lean`. The Lean declarations must preserve the physical quantities, dimensions or dimensional roles, figure labels, governing-law hypotheses, and final relation expressed by this problem.
Use Mathlib/Physlib names found through LeanExplore where available. If a physics API is missing, introduce faithful local abstractions rather than scalar placeholder aliases.

\begin{theorem}[Physics formalization target]
\label{thm:physics:phyx_mini_0775:target}
\lean{PhyXMiniProblems.ProblemPhyXMini0775.problem_phyx_mini_0775}
\uses{def:physics:phyx-mini-0775:phyxminiproblems-problemphyxmini0775-speedinmeterspersecond, def:physics:phyx-mini-0775:phyxminiproblems-problemphyxmini0775-weldeddiskpulleysetup, def:physics:phyx-mini-0775:phyxminiproblems-problemphyxmini0775-matchesprimaryfigure, def:physics:phyx-mini-0775:phyxminiproblems-problemphyxmini0775-matchesproblemdata, def:physics:phyx-mini-0775:phyxminiproblems-problemphyxmini0775-usesstandardearthgravity, def:physics:phyx-mini-0775:phyxminiproblems-problemphyxmini0775-hasphysicalparameters, def:physics:phyx-mini-0775:phyxminiproblems-problemphyxmini0775-satisfiesweldeduniformdiskrotationallaws, def:physics:phyx-mini-0775:phyxminiproblems-problemphyxmini0775-satisfiesnoslipstringkinematics, def:physics:phyx-mini-0775:phyxminiproblems-problemphyxmini0775-satisfiesreleasetoimpactenergyconservation, def:physics:phyx-mini-0775:phyxminiproblems-problemphyxmini0775-recordeddatasetanswer, def:physics:phyx-mini-0775:phyxminiproblems-problemphyxmini0775-roundstodisplayedspeed}
The assigned autoformalize agent should translate this physics problem into Lean declarations in the covered file, with theorem and lemma proof bodies written as `by sorry`.
\end{theorem}
\begin{proof}
This is an autoformalization task, not a proof task. Produce statements that can later be proved by the physics prover without weakening the physical model.
\end{proof}

% --- Archon declaration topology begin ---
\section{Declaration topology}
\begin{definition}[moment Of Inertia Dimension]
  \label{def:physics:phyx-mini-0775:phyxminiproblems-problemphyxmini0775-momentofinertiadimension}
  \lean{PhyXMiniProblems.ProblemPhyXMini0775.momentOfInertiaDimension}
  The physical dimension of an axial moment of inertia, mass * length²
\end{definition}
\begin{proof}
  This is the declared model definition of moment Of Inertia Dimension.
\end{proof}
\begin{definition}[Mass Quantity]
  \label{def:physics:phyx-mini-0775:phyxminiproblems-problemphyxmini0775-massquantity}
  \lean{PhyXMiniProblems.ProblemPhyXMini0775.MassQuantity}
  A nonnegative, unit-independent physical mass
\end{definition}
\begin{proof}
  This is the declared model definition of Mass Quantity.
\end{proof}
\begin{definition}[Length Quantity]
  \label{def:physics:phyx-mini-0775:phyxminiproblems-problemphyxmini0775-lengthquantity}
  \lean{PhyXMiniProblems.ProblemPhyXMini0775.LengthQuantity}
  A nonnegative, unit-independent physical length
\end{definition}
\begin{proof}
  This is the declared model definition of Length Quantity.
\end{proof}
\begin{definition}[Speed Quantity]
  \label{def:physics:phyx-mini-0775:phyxminiproblems-problemphyxmini0775-speedquantity}
  \lean{PhyXMiniProblems.ProblemPhyXMini0775.SpeedQuantity}
  A nonnegative physical speed
\end{definition}
\begin{proof}
  This is the declared model definition of Speed Quantity.
\end{proof}
\begin{definition}[Acceleration Quantity]
  \label{def:physics:phyx-mini-0775:phyxminiproblems-problemphyxmini0775-accelerationquantity}
  \lean{PhyXMiniProblems.ProblemPhyXMini0775.AccelerationQuantity}
  A nonnegative acceleration magnitude
\end{definition}
\begin{proof}
  This is the declared model definition of Acceleration Quantity.
\end{proof}
\begin{definition}[Angular Speed Quantity]
  \label{def:physics:phyx-mini-0775:phyxminiproblems-problemphyxmini0775-angularspeedquantity}
  \lean{PhyXMiniProblems.ProblemPhyXMini0775.AngularSpeedQuantity}
  A nonnegative angular-speed magnitude; radians are dimensionless
\end{definition}
\begin{proof}
  This is the declared model definition of Angular Speed Quantity.
\end{proof}
\begin{definition}[Moment Of Inertia Quantity]
  \label{def:physics:phyx-mini-0775:phyxminiproblems-problemphyxmini0775-momentofinertiaquantity}
  \lean{PhyXMiniProblems.ProblemPhyXMini0775.MomentOfInertiaQuantity}
  \uses{def:physics:phyx-mini-0775:phyxminiproblems-problemphyxmini0775-momentofinertiadimension}
  A nonnegative axial moment of inertia
\end{definition}
\begin{proof}
  This is the declared model definition of Moment Of Inertia Quantity.
\end{proof}
\begin{definition}[nonnegative SIReadout]
  \label{def:physics:phyx-mini-0775:phyxminiproblems-problemphyxmini0775-nonnegativesireadout}
  \lean{PhyXMiniProblems.ProblemPhyXMini0775.nonnegativeSIReadout}
  Read a nonnegative dimensionful quantity in coherent SI units
\end{definition}
\begin{proof}
  This is the declared model definition of nonnegative SIReadout.
\end{proof}
\begin{definition}[mass In Kilograms]
  \label{def:physics:phyx-mini-0775:phyxminiproblems-problemphyxmini0775-massinkilograms}
  \lean{PhyXMiniProblems.ProblemPhyXMini0775.massInKilograms}
  \uses{def:physics:phyx-mini-0775:phyxminiproblems-problemphyxmini0775-massquantity, def:physics:phyx-mini-0775:phyxminiproblems-problemphyxmini0775-nonnegativesireadout}
  Kilogram readout of a physical mass
\end{definition}
\begin{proof}
  This is the declared model definition of mass In Kilograms.
\end{proof}
\begin{definition}[length In Meters]
  \label{def:physics:phyx-mini-0775:phyxminiproblems-problemphyxmini0775-lengthinmeters}
  \lean{PhyXMiniProblems.ProblemPhyXMini0775.lengthInMeters}
  \uses{def:physics:phyx-mini-0775:phyxminiproblems-problemphyxmini0775-lengthquantity, def:physics:phyx-mini-0775:phyxminiproblems-problemphyxmini0775-nonnegativesireadout}
  Meter readout of a physical length
\end{definition}
\begin{proof}
  This is the declared model definition of length In Meters.
\end{proof}
\begin{definition}[speed In Meters Per Second]
  \label{def:physics:phyx-mini-0775:phyxminiproblems-problemphyxmini0775-speedinmeterspersecond}
  \lean{PhyXMiniProblems.ProblemPhyXMini0775.speedInMetersPerSecond}
  \uses{def:physics:phyx-mini-0775:phyxminiproblems-problemphyxmini0775-speedquantity}
  Meter-per-second readout of a physical speed
\end{definition}
\begin{proof}
  This is the declared model definition of speed In Meters Per Second.
\end{proof}
\begin{definition}[acceleration In Meters Per Second Squared]
  \label{def:physics:phyx-mini-0775:phyxminiproblems-problemphyxmini0775-accelerationinmeterspersecondsquared}
  \lean{PhyXMiniProblems.ProblemPhyXMini0775.accelerationInMetersPerSecondSquared}
  \uses{def:physics:phyx-mini-0775:phyxminiproblems-problemphyxmini0775-accelerationquantity, def:physics:phyx-mini-0775:phyxminiproblems-problemphyxmini0775-nonnegativesireadout}
  Meter-per-second-squared readout of an acceleration magnitude
\end{definition}
\begin{proof}
  This is the declared model definition of acceleration In Meters Per Second Squared.
\end{proof}
\begin{definition}[angular Speed In Radians Per Second]
  \label{def:physics:phyx-mini-0775:phyxminiproblems-problemphyxmini0775-angularspeedinradianspersecond}
  \lean{PhyXMiniProblems.ProblemPhyXMini0775.angularSpeedInRadiansPerSecond}
  \uses{def:physics:phyx-mini-0775:phyxminiproblems-problemphyxmini0775-angularspeedquantity, def:physics:phyx-mini-0775:phyxminiproblems-problemphyxmini0775-nonnegativesireadout}
  Radian-per-second readout of an angular-speed magnitude
\end{definition}
\begin{proof}
  This is the declared model definition of angular Speed In Radians Per Second.
\end{proof}
\begin{definition}[moment Of Inertia In Kilogram Meters Squared]
  \label{def:physics:phyx-mini-0775:phyxminiproblems-problemphyxmini0775-momentofinertiainkilogrammeterssquared}
  \lean{PhyXMiniProblems.ProblemPhyXMini0775.momentOfInertiaInKilogramMetersSquared}
  \uses{def:physics:phyx-mini-0775:phyxminiproblems-problemphyxmini0775-momentofinertiaquantity, def:physics:phyx-mini-0775:phyxminiproblems-problemphyxmini0775-nonnegativesireadout}
  Kilogram-meter-squared readout of an axial moment of inertia
\end{definition}
\begin{proof}
  This is the declared model definition of moment Of Inertia In Kilogram Meters Squared.
\end{proof}
\begin{definition}[Disk]
  \label{def:physics:phyx-mini-0775:phyxminiproblems-problemphyxmini0775-disk}
  \lean{PhyXMiniProblems.ProblemPhyXMini0775.Disk}
  The two welded disks
\end{definition}
\begin{proof}
  This is the declared model definition of Disk.
\end{proof}
\begin{definition}[Motion Stage]
  \label{def:physics:phyx-mini-0775:phyxminiproblems-problemphyxmini0775-motionstage}
  \lean{PhyXMiniProblems.ProblemPhyXMini0775.MotionStage}
  The two stages needed for the release-to-impact energy balance
\end{definition}
\begin{proof}
  This is the declared model definition of Motion Stage.
\end{proof}
\begin{definition}[Radius Label]
  \label{def:physics:phyx-mini-0775:phyxminiproblems-problemphyxmini0775-radiuslabel}
  \lean{PhyXMiniProblems.ProblemPhyXMini0775.RadiusLabel}
  Radius labels visible in the primary image
\end{definition}
\begin{proof}
  This is the declared model definition of Radius Label.
\end{proof}
\begin{definition}[Disk Mass Model]
  \label{def:physics:phyx-mini-0775:phyxminiproblems-problemphyxmini0775-diskmassmodel}
  \lean{PhyXMiniProblems.ProblemPhyXMini0775.DiskMassModel}
  The mass-distribution idealization used by the disk inertia law
\end{definition}
\begin{proof}
  This is the declared model definition of Disk Mass Model.
\end{proof}
\begin{definition}[Disk Coupling]
  \label{def:physics:phyx-mini-0775:phyxminiproblems-problemphyxmini0775-diskcoupling}
  \lean{PhyXMiniProblems.ProblemPhyXMini0775.DiskCoupling}
  Mechanical coupling between the two concentric disks
\end{definition}
\begin{proof}
  This is the declared model definition of Disk Coupling.
\end{proof}
\begin{definition}[Axle Location]
  \label{def:physics:phyx-mini-0775:phyxminiproblems-problemphyxmini0775-axlelocation}
  \lean{PhyXMiniProblems.ProblemPhyXMini0775.AxleLocation}
  Location of the axle relative to the disks
\end{definition}
\begin{proof}
  This is the declared model definition of Axle Location.
\end{proof}
\begin{definition}[Axle Orientation]
  \label{def:physics:phyx-mini-0775:phyxminiproblems-problemphyxmini0775-axleorientation}
  \lean{PhyXMiniProblems.ProblemPhyXMini0775.AxleOrientation}
  Orientation of the axle in the supplied image
\end{definition}
\begin{proof}
  This is the declared model definition of Axle Orientation.
\end{proof}
\begin{definition}[Axle Resistance Model]
  \label{def:physics:phyx-mini-0775:phyxminiproblems-problemphyxmini0775-axleresistancemodel}
  \lean{PhyXMiniProblems.ProblemPhyXMini0775.AxleResistanceModel}
  Resistive idealization of the rotation axis
\end{definition}
\begin{proof}
  This is the declared model definition of Axle Resistance Model.
\end{proof}
\begin{definition}[String Mass Model]
  \label{def:physics:phyx-mini-0775:phyxminiproblems-problemphyxmini0775-stringmassmodel}
  \lean{PhyXMiniProblems.ProblemPhyXMini0775.StringMassModel}
  Mass idealization justified by the description of the string as light
\end{definition}
\begin{proof}
  This is the declared model definition of String Mass Model.
\end{proof}
\begin{definition}[String Disk Contact]
  \label{def:physics:phyx-mini-0775:phyxminiproblems-problemphyxmini0775-stringdiskcontact}
  \lean{PhyXMiniProblems.ProblemPhyXMini0775.StringDiskContact}
  Contact condition between the string and the smaller disk
\end{definition}
\begin{proof}
  This is the declared model definition of String Disk Contact.
\end{proof}
\begin{definition}[Welded Disk Pulley Figure]
  \label{def:physics:phyx-mini-0775:phyxminiproblems-problemphyxmini0775-weldeddiskpulleyfigure}
  \lean{PhyXMiniProblems.ProblemPhyXMini0775.WeldedDiskPulleyFigure}
  \uses{def:physics:phyx-mini-0775:phyxminiproblems-problemphyxmini0775-disk, def:physics:phyx-mini-0775:phyxminiproblems-problemphyxmini0775-radiuslabel}
  Qualitative evidence transcribed from image 775.png. The image supplies layout and labels only; no physical radius is inferred from the drawn sizes
\end{definition}
\begin{proof}
  This is the declared model definition of Welded Disk Pulley Figure.
\end{proof}
\begin{definition}[Welded Disk Pulley Setup]
  \label{def:physics:phyx-mini-0775:phyxminiproblems-problemphyxmini0775-weldeddiskpulleysetup}
  \lean{PhyXMiniProblems.ProblemPhyXMini0775.WeldedDiskPulleySetup}
  \uses{def:physics:phyx-mini-0775:phyxminiproblems-problemphyxmini0775-massquantity, def:physics:phyx-mini-0775:phyxminiproblems-problemphyxmini0775-lengthquantity, def:physics:phyx-mini-0775:phyxminiproblems-problemphyxmini0775-speedquantity, def:physics:phyx-mini-0775:phyxminiproblems-problemphyxmini0775-accelerationquantity, def:physics:phyx-mini-0775:phyxminiproblems-problemphyxmini0775-angularspeedquantity, def:physics:phyx-mini-0775:phyxminiproblems-problemphyxmini0775-momentofinertiaquantity, def:physics:phyx-mini-0775:phyxminiproblems-problemphyxmini0775-disk, def:physics:phyx-mini-0775:phyxminiproblems-problemphyxmini0775-motionstage, def:physics:phyx-mini-0775:phyxminiproblems-problemphyxmini0775-diskmassmodel, def:physics:phyx-mini-0775:phyxminiproblems-problemphyxmini0775-diskcoupling, def:physics:phyx-mini-0775:phyxminiproblems-problemphyxmini0775-axlelocation, def:physics:phyx-mini-0775:phyxminiproblems-problemphyxmini0775-axleorientation, def:physics:phyx-mini-0775:phyxminiproblems-problemphyxmini0775-axleresistancemodel, def:physics:phyx-mini-0775:phyxminiproblems-problemphyxmini0775-stringmassmodel, def:physics:phyx-mini-0775:phyxminiproblems-problemphyxmini0775-stringdiskcontact, def:physics:phyx-mini-0775:phyxminiproblems-problemphyxmini0775-weldeddiskpulleyfigure}
  Independent physical quantities for the system. Speeds at both stages, angular speeds, component and total inertias, and the rigid-body realization are stored independently. None is defined from the displayed 3.40 m/s answer
\end{definition}
\begin{proof}
  This is the declared model definition of Welded Disk Pulley Setup.
\end{proof}
\begin{definition}[Matches Primary Figure]
  \label{def:physics:phyx-mini-0775:phyxminiproblems-problemphyxmini0775-matchesprimaryfigure}
  \lean{PhyXMiniProblems.ProblemPhyXMini0775.MatchesPrimaryFigure}
  \uses{def:physics:phyx-mini-0775:phyxminiproblems-problemphyxmini0775-weldeddiskpulleysetup}
  Facts visible in the supplied primary image
\end{definition}
\begin{proof}
  This is the declared model definition of Matches Primary Figure.
\end{proof}
\begin{definition}[Matches Problem Data]
  \label{def:physics:phyx-mini-0775:phyxminiproblems-problemphyxmini0775-matchesproblemdata}
  \lean{PhyXMiniProblems.ProblemPhyXMini0775.MatchesProblemData}
  \uses{def:physics:phyx-mini-0775:phyxminiproblems-problemphyxmini0775-massinkilograms, def:physics:phyx-mini-0775:phyxminiproblems-problemphyxmini0775-lengthinmeters, def:physics:phyx-mini-0775:phyxminiproblems-problemphyxmini0775-speedinmeterspersecond, def:physics:phyx-mini-0775:phyxminiproblems-problemphyxmini0775-angularspeedinradianspersecond, def:physics:phyx-mini-0775:phyxminiproblems-problemphyxmini0775-weldeddiskpulleysetup}
  Numerical and qualitative data supplied in the prose. Uniform solid disks and a negligible-mass taut string are the standard textbook idealizations needed for the recorded numerical answer
\end{definition}
\begin{proof}
  This is the declared model definition of Matches Problem Data.
\end{proof}
\begin{definition}[Uses Standard Earth Gravity]
  \label{def:physics:phyx-mini-0775:phyxminiproblems-problemphyxmini0775-usesstandardearthgravity}
  \lean{PhyXMiniProblems.ProblemPhyXMini0775.UsesStandardEarthGravity}
  \uses{def:physics:phyx-mini-0775:phyxminiproblems-problemphyxmini0775-accelerationinmeterspersecondsquared, def:physics:phyx-mini-0775:phyxminiproblems-problemphyxmini0775-weldeddiskpulleysetup}
  The standard near-Earth value used to evaluate the displayed choices
\end{definition}
\begin{proof}
  This is the declared model definition of Uses Standard Earth Gravity.
\end{proof}
\begin{definition}[Has Physical Parameters]
  \label{def:physics:phyx-mini-0775:phyxminiproblems-problemphyxmini0775-hasphysicalparameters}
  \lean{PhyXMiniProblems.ProblemPhyXMini0775.HasPhysicalParameters}
  \uses{def:physics:phyx-mini-0775:phyxminiproblems-problemphyxmini0775-massinkilograms, def:physics:phyx-mini-0775:phyxminiproblems-problemphyxmini0775-lengthinmeters, def:physics:phyx-mini-0775:phyxminiproblems-problemphyxmini0775-accelerationinmeterspersecondsquared, def:physics:phyx-mini-0775:phyxminiproblems-problemphyxmini0775-momentofinertiainkilogrammeterssquared, def:physics:phyx-mini-0775:phyxminiproblems-problemphyxmini0775-weldeddiskpulleysetup}
  Positivity and nondegeneracy of the physical inputs. In particular, no numerical value or positivity beyond nonnegativity is imposed on the unknown impact speed
\end{definition}
\begin{proof}
  This is the declared model definition of Has Physical Parameters.
\end{proof}
\begin{definition}[Satisfies Welded Uniform Disk Rotational Laws]
  \label{def:physics:phyx-mini-0775:phyxminiproblems-problemphyxmini0775-satisfiesweldeduniformdiskrotationallaws}
  \lean{PhyXMiniProblems.ProblemPhyXMini0775.SatisfiesWeldedUniformDiskRotationalLaws}
  \uses{def:physics:phyx-mini-0775:phyxminiproblems-problemphyxmini0775-massinkilograms, def:physics:phyx-mini-0775:phyxminiproblems-problemphyxmini0775-lengthinmeters, def:physics:phyx-mini-0775:phyxminiproblems-problemphyxmini0775-angularspeedinradianspersecond, def:physics:phyx-mini-0775:phyxminiproblems-problemphyxmini0775-momentofinertiainkilogrammeterssquared, def:physics:phyx-mini-0775:phyxminiproblems-problemphyxmini0775-weldeddiskpulleysetup}
  The two uniform-solid-disk inertia formulas, their additivity after welding, and the bridge to Physlib's three-dimensional rigid-body representation. Angular velocity is directed along the selected axle coordinate
\end{definition}
\begin{proof}
  This is the declared model definition of Satisfies Welded Uniform Disk Rotational Laws.
\end{proof}
\begin{definition}[Satisfies No Slip String Kinematics]
  \label{def:physics:phyx-mini-0775:phyxminiproblems-problemphyxmini0775-satisfiesnoslipstringkinematics}
  \lean{PhyXMiniProblems.ProblemPhyXMini0775.SatisfiesNoSlipStringKinematics}
  \uses{def:physics:phyx-mini-0775:phyxminiproblems-problemphyxmini0775-lengthinmeters, def:physics:phyx-mini-0775:phyxminiproblems-problemphyxmini0775-speedinmeterspersecond, def:physics:phyx-mini-0775:phyxminiproblems-problemphyxmini0775-angularspeedinradianspersecond, def:physics:phyx-mini-0775:phyxminiproblems-problemphyxmini0775-weldeddiskpulleysetup}
  A taut string that does not slip has block speed equal to the tangential speed of the smaller disk, v = R₁ * omega, at each modeled stage
\end{definition}
\begin{proof}
  This is the declared model definition of Satisfies No Slip String Kinematics.
\end{proof}
\begin{definition}[Satisfies Release To Impact Energy Conservation]
  \label{def:physics:phyx-mini-0775:phyxminiproblems-problemphyxmini0775-satisfiesreleasetoimpactenergyconservation}
  \lean{PhyXMiniProblems.ProblemPhyXMini0775.SatisfiesReleaseToImpactEnergyConservation}
  \uses{def:physics:phyx-mini-0775:phyxminiproblems-problemphyxmini0775-massinkilograms, def:physics:phyx-mini-0775:phyxminiproblems-problemphyxmini0775-lengthinmeters, def:physics:phyx-mini-0775:phyxminiproblems-problemphyxmini0775-speedinmeterspersecond, def:physics:phyx-mini-0775:phyxminiproblems-problemphyxmini0775-accelerationinmeterspersecondsquared, def:physics:phyx-mini-0775:phyxminiproblems-problemphyxmini0775-weldeddiskpulleysetup}
  Mechanical energy conservation between release and the instant immediately before impact. The three terms are block gravitational potential energy, block translational kinetic energy, and the welded disks' rotational kinetic energy from RigidBody.rotationalKineticEnergy. This is a governing relation. It contains the independent impact speed but does not state its solved value or mention an answer choice
\end{definition}
\begin{proof}
  This is the declared model definition of Satisfies Release To Impact Energy Conservation.
\end{proof}
\begin{lemma}[total Moment Of Inertia eq nine over four thousand]
  \label{lem:physics:phyx-mini-0775:phyxminiproblems-problemphyxmini0775-totalmomentofinertia-eq-nine-over-four-thousand}
  \lean{PhyXMiniProblems.ProblemPhyXMini0775.totalMomentOfInertia_eq_nine_over_four_thousand}
  \uses{def:physics:phyx-mini-0775:phyxminiproblems-problemphyxmini0775-momentofinertiainkilogrammeterssquared, def:physics:phyx-mini-0775:phyxminiproblems-problemphyxmini0775-weldeddiskpulleysetup, def:physics:phyx-mini-0775:phyxminiproblems-problemphyxmini0775-matchesproblemdata, def:physics:phyx-mini-0775:phyxminiproblems-problemphyxmini0775-satisfiesweldeduniformdiskrotationallaws}
  The two component inertias add to 0.00225 kg m² = 9/4000 kg m²
\end{lemma}
\begin{proof}
  The conclusion follows from the governing relations and figure readouts named in the statement of total Moment Of Inertia eq nine over four thousand.
\end{proof}
\begin{lemma}[impact Speed squared formula]
  \label{lem:physics:phyx-mini-0775:phyxminiproblems-problemphyxmini0775-impactspeed-squared-formula}
  \lean{PhyXMiniProblems.ProblemPhyXMini0775.impactSpeed_squared_formula}
  \uses{def:physics:phyx-mini-0775:phyxminiproblems-problemphyxmini0775-massinkilograms, def:physics:phyx-mini-0775:phyxminiproblems-problemphyxmini0775-lengthinmeters, def:physics:phyx-mini-0775:phyxminiproblems-problemphyxmini0775-speedinmeterspersecond, def:physics:phyx-mini-0775:phyxminiproblems-problemphyxmini0775-accelerationinmeterspersecondsquared, def:physics:phyx-mini-0775:phyxminiproblems-problemphyxmini0775-momentofinertiainkilogrammeterssquared, def:physics:phyx-mini-0775:phyxminiproblems-problemphyxmini0775-weldeddiskpulleysetup, def:physics:phyx-mini-0775:phyxminiproblems-problemphyxmini0775-matchesproblemdata, def:physics:phyx-mini-0775:phyxminiproblems-problemphyxmini0775-hasphysicalparameters, def:physics:phyx-mini-0775:phyxminiproblems-problemphyxmini0775-satisfiesweldeduniformdiskrotationallaws, def:physics:phyx-mini-0775:phyxminiproblems-problemphyxmini0775-satisfiesnoslipstringkinematics, def:physics:phyx-mini-0775:phyxminiproblems-problemphyxmini0775-satisfiesreleasetoimpactenergyconservation}
  Eliminating angular speed with v = R₁ * omega gives the standard effective- mass formula. It is a consequence of the governing laws, not a premise
\end{lemma}
\begin{proof}
  The conclusion follows from the governing relations and figure readouts named in the statement of impact Speed squared formula.
\end{proof}
\begin{lemma}[impact Speed eq sqrt 196 over 17]
  \label{lem:physics:phyx-mini-0775:phyxminiproblems-problemphyxmini0775-impactspeed-eq-sqrt-196-over-17}
  \lean{PhyXMiniProblems.ProblemPhyXMini0775.impactSpeed_eq_sqrt_196_over_17}
  \uses{def:physics:phyx-mini-0775:phyxminiproblems-problemphyxmini0775-speedinmeterspersecond, def:physics:phyx-mini-0775:phyxminiproblems-problemphyxmini0775-weldeddiskpulleysetup, def:physics:phyx-mini-0775:phyxminiproblems-problemphyxmini0775-matchesproblemdata, def:physics:phyx-mini-0775:phyxminiproblems-problemphyxmini0775-usesstandardearthgravity, def:physics:phyx-mini-0775:phyxminiproblems-problemphyxmini0775-hasphysicalparameters, def:physics:phyx-mini-0775:phyxminiproblems-problemphyxmini0775-satisfiesweldeduniformdiskrotationallaws, def:physics:phyx-mini-0775:phyxminiproblems-problemphyxmini0775-satisfiesnoslipstringkinematics, def:physics:phyx-mini-0775:phyxminiproblems-problemphyxmini0775-satisfiesreleasetoimpactenergyconservation}
  For the supplied disk and block data with g = 9.8 m/s², the nonnegative impact speed is exactly sqrt (196/17) m/s, approximately 3.3955 m/s
\end{lemma}
\begin{proof}
  The conclusion follows from the governing relations and figure readouts named in the statement of impact Speed eq sqrt 196 over 17.
\end{proof}
\begin{definition}[Answer Choice]
  \label{def:physics:phyx-mini-0775:phyxminiproblems-problemphyxmini0775-answerchoice}
  \lean{PhyXMiniProblems.ProblemPhyXMini0775.AnswerChoice}
  Labels of the four speeds displayed with the exercise
\end{definition}
\begin{proof}
  This is the declared model definition of Answer Choice.
\end{proof}
\begin{definition}[speed In Meters Per Second]
  \label{def:physics:phyx-mini-0775:phyxminiproblems-problemphyxmini0775-answerchoice-speedinmeterspersecond}
  \lean{PhyXMiniProblems.ProblemPhyXMini0775.AnswerChoice.speedInMetersPerSecond}
  \uses{def:physics:phyx-mini-0775:phyxminiproblems-problemphyxmini0775-speedinmeterspersecond, def:physics:phyx-mini-0775:phyxminiproblems-problemphyxmini0775-answerchoice}
  Meter-per-second value printed beside each answer label
\end{definition}
\begin{proof}
  This is the declared model definition of speed In Meters Per Second.
\end{proof}
\begin{definition}[recorded Dataset Answer]
  \label{def:physics:phyx-mini-0775:phyxminiproblems-problemphyxmini0775-recordeddatasetanswer}
  \lean{PhyXMiniProblems.ProblemPhyXMini0775.recordedDatasetAnswer}
  \uses{def:physics:phyx-mini-0775:phyxminiproblems-problemphyxmini0775-answerchoice}
  The answer label recorded in the supplied dataset metadata
\end{definition}
\begin{proof}
  This is the declared model definition of recorded Dataset Answer.
\end{proof}
\begin{definition}[Rounds To Displayed Speed]
  \label{def:physics:phyx-mini-0775:phyxminiproblems-problemphyxmini0775-roundstodisplayedspeed}
  \lean{PhyXMiniProblems.ProblemPhyXMini0775.RoundsToDisplayedSpeed}
  \uses{def:physics:phyx-mini-0775:phyxminiproblems-problemphyxmini0775-speedquantity, def:physics:phyx-mini-0775:phyxminiproblems-problemphyxmini0775-speedinmeterspersecond, def:physics:phyx-mini-0775:phyxminiproblems-problemphyxmini0775-answerchoice}
  Agreement with a displayed hundredth means lying in its half-cent-wide rounding interval. This avoids the false statement that the unrounded exact speed is literally 3.40 m/s
\end{definition}
\begin{proof}
  This is the declared model definition of Rounds To Displayed Speed.
\end{proof}
% --- Archon declaration topology end ---
% --- Archon physics formalization source end ---
